\section{The corrected identity and a prime-carrying smooth transport}
\label{sec:trace-global}

\subsection{The assertion, the actual objects, and the zero-location ledger}

Throughout this section, a statement attributed to a predecessor is used
with its recorded hypotheses. In particular, a theorem about convolution
on a quotient is not silently changed into a theorem about an arbitrary
lift of that convolution to its ambient Hilbert space.

Write $C_{\mathbb Q}=\widehat{\mathbb Z}^{\times}\times\mathbb R_{>0}$,
with probability Haar measure on the first factor and $dt/t$ on the
second. Put $x=\log t$, and let $\mathcal E$ be the space of smooth
functions on $\mathbb R$ all of whose derivatives decay faster than
$e^{-N|x|}$ for every $N$. Its Mellin transform is
\[
 \mathcal Mf(z)=\int_{\mathbb R}f(x)e^{zx}\,dx.
\]
For the trivial character let $\mathcal I$ be the closed convolution
ideal of GRH\_NIGHT MT-M06/MT-M18. MT-M24 identifies it, without
zero-location input, with the actual theta-summation image. The
centering operator is $T_{1/2}f(x)=e^{x/2}f(x)$. The corrected strip
space is
\[
 H_{\alpha,s}=L^2\bigl(\mathbb R,
      e^{2\alpha|x|}(1+x^2)^{s/2}\,dx\bigr),\qquad
 J_{\alpha,s}=\overline{T_{1/2}\mathcal I}^{\,H_{\alpha,s}},
 \qquad Z_{\alpha,s}=H_{\alpha,s}/J_{\alpha,s}.
 \tag{MA-TR-1}
\]
Here $\alpha>0$; a fixed $s\ge3$ is compatible with the source
geometric-kernel convention. Let $P_G$ and $P_Z$ be the orthogonal
projections onto $J_{\alpha,s}$ and its orthogonal complement. These
are Masullo's Definition 10.4.1.1, p.~456, on the corrected closed
subspace. For $h\in\mathcal D^0:=\{h\in C_c^\infty(\mathbb R_{>0}):
h(1)=0\}$, set
\[
 \widehat h(z)=\int_0^\infty h(t)t^z\frac{dt}{t},\qquad
 C_hf(x)=\int_{\mathbb R}h(e^y)f(x-y)\,dy.
\]
The source's convolution-after-projection operators are
$S_h=C_hP_Z$ and $G_{\rho,h}=M_\rho C_hP_G$; see Definition
10.4.1.5, p.~458, and (10.4.3.45), (10.4.3.48), p.~473.
The other order, $P_ZC_h$, is explicitly distinguished in
(10.4.1.10)--(10.4.1.11), pp.~457--458.

The corrected indentation may be taken to be the indicator of
$\mathbb R_{>0}\times\widehat{\mathbb Z}^{\times}$ inside the ideles,
as in ALPOGE\_BRIDGES F-SM-20. Its compact compression is valid, but
does not identify the compressed global operator with a separately
constructed local one. The closure and cutoff corrections are precisely
GRH\_NIGHT MT-M31--MT-M35; no injectivity of an indicator multiplier
or interchange of a cutoff image with a cofinal intersection is used.

In these normalizations put
\[
\begin{split}
 P(h)&=\widehat h(1/2)+\widehat h(-1/2),\\
 A_0(h)&=\int_0^\infty
   \frac{e^{-x/2}}{1-e^{-2x}}
       \bigl(h(e^x)+h(e^{-x})\bigr)\,dx,\\
 B(h)&=\sum_p\sum_{k\ge1}\frac{\log p}{p^{k/2}}
                   \bigl(h(p^k)+h(p^{-k})\bigr),\\
 \tau_\alpha(h)&=\sum_{|\Re\rho-1/2|<\alpha}
                         m_\rho\widehat h(\rho-1/2).
\end{split}
\tag{MA-TR-2}
\]
All the displayed integrals and sums converge absolutely. The prime
sum is finite for each compact test. Smoothness and $h(1)=0$ remove
the archimedean singularity. The spectral convergence uses only the
classical zero count and rapid vertical decay of $\widehat h$.

There is an essential distinction between the open-strip distribution
$\tau_\alpha$ just defined and an already established trace on
$Z_{\alpha,s}$. MT-M26 identifies the actual quotient trace when
\emph{all} zeros lie in a closed substrip of its interior. This holds
unconditionally for $\alpha>1/2$, and for every $\alpha>0$ under
$A$ in the present rational case. It does not, as stated, prove an
unconditional trace-class theorem for a partial strip. That additional theorem is proved later in MA-WP-001--003,
with its full boundary multiplicity convention; it is not
imported from MT-M26. At a boundary zero the polynomial
weight also retains some continuous jet functionals, as MA-TR-103
below makes precise. Thus a literal identification with an open-strip
sum would require a boundary convention or a boundary-free width.

The requested central geometric identity is
\[
 \boxed{\displaystyle
  \lim_{\rho\to1}\operatorname{Tr}G_{\rho,h}
     = A_0(h)+B(h)-P(h)=-\tau_\alpha(h)
       \quad(\alpha>0,\ h\in\mathcal D^0).}
 \tag{MA-TR-3}
\]
The first equality is the geometric transport assertion. The last is
the narrow spectral identity, with the source's sign convention.
Where the actual Hilbert trace is invoked, the preceding scope and
boundary qualification are part of the assertion. Proving it without
zero-location input includes, rather than presupposes, those missing
operator and comparison statements.
The local theorem alone gives $A_0(h)$ for $S=\{\infty\}$, before
the two removed moments, and gives $A_0(h)$ plus the primes in $S$
for larger $S$. It does not give the missing identification of global
operators. The exact final source location is the use of Theorem
12.4.1.1 in the proof on p.~520, followed by (12.5.1.4). Its proposed
bridge is Section~12.3, p.~514. The latter contains no proof.

The classical full explicit formula, in the conventions already
checked by F-SM-13 and MT-M30, is
\[
 A_0(h)+B(h)=P(h)-\tau_{\rm full}(h).
 \tag{MA-TR-4}
\]
Consequently the exact unconditional discrepancy between the numerical
geometric expression and the proposed restricted spectral expression is
\[
 A_0(h)+B(h)-P(h)+\tau_\alpha(h)
   =-\!\sum_{|\Re\rho-1/2|\ge\alpha}
                    m_\rho\widehat h(\rho-1/2).
 \tag{MA-TR-5}
\]
The equality criterion for these restricted tests is already proved in
MT-M30, (M30.14); it is not reproved or claimed as new here. Theorem
$A$ makes (MA-TR-5) zero for the rational and Dirichlet cases. Using
that fact in a proposed independent proof of (MA-TR-3) would be
circular with respect to the method being tested, even though it is a
valid application of an external theorem.

\begin{center}
\begin{tabular}{p{0.51\linewidth}p{0.40\linewidth}}
\hline
Step & Zero-location input \\
\hline
Theta image, closed ideal, compact averaging and valid indentation &
None; MT-M18/M24 and F-SM-20.\\
Local Euler factors, the corrected Poisson calculation, and the local
Fourier cutoff calculation & None.\\
Full spectral trace on a strip containing $[0,1]$ in its interior &
Classical containing strip only; MT-M26/M27.\\
Faithful spectral trace on every narrow strip & $A$ in the Dirichlet
case; $M$ in the general Hecke case, by MT-M26/M28.\\
Replacing the full zero sum by every narrow sum & Exactly the
zero-location assertion; circular in an independent proof.\\
Global-to-local transport of the specified cutoff kernels &
Not supplied by any of the foregoing.\\
\hline
\end{tabular}
\end{center}

\paragraph{MA-TR-103 (the endpoint jet convention).}
For the exact weight in (MA-TR-1), the functional
$f\mapsto\int x^j f(x)e^{\lambda x}\,dx$, initially on
$C_c^\infty(\mathbb R)$, extends continuously to $H_{\alpha,s}$
if and only if
\[
 |\Re\lambda|<\alpha,
 \quad\hbox{or}\quad
 |\Re\lambda|=\alpha\ \hbox{ and }\ 2j<s-1.
 \tag{MA-TR-12}
\]
It never extends when $|\Re\lambda|>\alpha$.
\emph{Proof.} By the weighted $L^2$ duality theorem, extension
is equivalent to finiteness of
\[
 \int_{\mathbb R}|x|^{2j}
        e^{2\Re\lambda x-2\alpha|x|}(1+x^2)^{-s/2}\,dx.
\]
The interior and exterior cases follow from exponential decay
and growth, respectively. At an endpoint one tail is comparable
to $\int_1^\infty x^{2j-s}\,dx$, which is finite precisely
under the displayed strict inequality. Density of compact smooth
functions proves necessity as well as sufficiency.
\hfill$\square$
This is the two-sided strip version of the dual-weight test in
GRH\_NIGHT MT-M25, not a proof of a complete partial-strip
spectral synthesis theorem. In particular, for $s\ge3$ every
boundary zero still imposes at least its value constraint on
$J_{\alpha,s}$.

\subsection{An explicit transport before the Hilbert completion}

\paragraph{MA-TR-101 (theta--M\"obius transport and the escape of its Euler tails).}
Let
\[
 \mathscr S_{0,+}=\{f\in\mathscr S(\mathbb R):
          f(-t)=f(t),\ f(0)=0,\ \int_{\mathbb R}f(t)\,dt=0\}.
\]
Associate to $f$ the radial function $F(t)=2\sqrt t\,f(t)$ for $t>0$.
The map
\[
 (\mathcal ZF)(t)=\sum_{n\ge1}n^{-1/2}F(nt)
                 =2\sqrt t\sum_{n\ge1}f(nt)
 \tag{MA-TR-6}
\]
is the centered global theta map in the trivial compact-character
sector. It is a continuous linear isomorphism from the displayed
Schwartz space onto $T_{1/2}\mathcal I$, with its smooth topology.
Its inverse on that image is explicitly
\[
 (\mathcal Jg)(t)=\sum_{n\ge1}\mu(n)n^{-1/2}g(nt).
 \tag{MA-TR-7}
\]
Both formulas converge with all derivatives locally uniformly for
$t>0$. They commute with norm translations on the $F$ variable
(the corresponding action on $f$ is
$f(t)\mapsto e^{-y/2}f(e^{-y}t)$). Thus there is a genuine
prime-carrying smooth transport; it has not yet been transported to
the source's narrow Hilbert kernels.

More precisely, for a finite set $S$ of rational primes define
\[
 P_S=\prod_{p\in S}(1-p^{-1/2}U_{-\log p}),\qquad
 (U_y g)(t)=g(e^{-y}t).
\]
Then
\[
 P_S\mathcal ZF(t)=
       \sum_{\substack{n\ge1\\(n,\prod_{p\in S}p)=1}}
                    n^{-1/2}F(nt).
 \tag{MA-TR-8}
\]
Suppose $F\in C_c^\infty([a,b])$, $0<a<b$, and
$\widehat F(1/2)=0$. Let $\rho$ be any nontrivial zeta zero with
$\lambda=\rho-1/2$ and $|\Re\lambda|<\alpha$. One may choose such
an $F$ with $\widehat F(\lambda)=1$. For this choice, every finite
$S$ satisfies the strict, $S$-independent bound
\[
 \left\|P_S\mathcal ZF-F\right\|_{H_{\alpha,s}}
 \ \ge\
 \left(\int_{\mathbb R}
     e^{2\Re\lambda x-2\alpha|x|}(1+x^2)^{-s/2}\,dx\right)^{-1/2}.
 \tag{MA-TR-9}
\]
If $S$ exhausts the primes and $q_S$ is the least omitted prime,
the difference in (MA-TR-9) is supported in $0<t\le b/q_S$ and
tends to zero in $C^\infty$ on every compact subset of
$\mathbb R_{>0}$. Thus local convergence of the prime-removal
operation conceals a nonvanishing Hilbert-space remainder escaping
to norm zero.

\emph{Proof.}
For fixed $t>0$, every derivative of the theta series converges
absolutely by the rapid decay of $f$ at infinity. Uniform convergence
holds when $t$ remains in a compact subinterval of $(0,\infty)$.
For large $t$ this also proves rapid decay of all logarithmic
derivatives. Poisson summation for the even Schwartz function gives
\[
 \sum_{n\in\mathbb Z}f(nt)=t^{-1}
                \sum_{n\in\mathbb Z}\widehat f(n/t).
\]
The two zero terms vanish by the conditions defining
$\mathscr S_{0,+}$. The right side and every logarithmic derivative
therefore decrease faster than every power of $t$ as $t\downarrow0$.
Multiplication by $\sqrt t$ preserves this conclusion. This proves
that $\mathcal ZF$ is exponential Schwartz in $\log t$, and the same
estimates, in terms of finitely many Schwartz seminorms for each
target seminorm, prove continuity.

For completeness, its image is the actual entire trivial-character
theta image, not just a collection of functions with the same
zeros. Average an adelic Schwartz representative over
$\widehat{\mathbb Z}^{\times}$. At a finite place, a locally constant
compactly supported unit-invariant function is a finite linear
combination of indicators of $p^j\mathbb Z_p$, $j\in\mathbb Z$.
Indeed its nonzero valuation shells are finite in number away from
a sufficiently small ball, on which local constancy makes it
constant; differences of adjacent ball indicators give the shells.
A restricted tensor product is therefore a finite linear combination
of rational dilates of $\prod_p1_{\mathbb Z_p}$. Absorb each rational
dilate into the real Schwartz function. A ball $q\widehat{\mathbb Z}$
contributes precisely the real function $f(qt)$ to the theta sum;
its finite additive volume is $q^{-1}$ for $q>0$. Thus both the
value at zero and the integral condition are preserved under this
absorption. Replacing the resulting real function by its even part
does not change its theta sum and preserves those conditions.
This proves surjectivity onto the actual compact-trivial image.
MT-M24 identifies that image with the closed ideal used here.

For $g=\mathcal ZF$ and $t>0$, the double series obtained from
(MA-TR-7) is absolutely convergent. For example, with $t$ in a fixed
compact interval and a sufficiently large $N$ its absolute value
and each logarithmic derivative are bounded by a constant times
\[
 \sum_{m,n\ge1}(mn)^{-1/2}(1+mnt)^{-N}<\infty.
\]
The coefficient of $k^{-1/2}F(kt)$ after regrouping is
$\sum_{n\mid k}\mu(n)$, which is $1$ at $k=1$ and $0$ otherwise.
This proves $\mathcal J\mathcal ZF=F$ and injectivity. The source
and target are Fr\'echet spaces: the source is a closed subspace
of Schwartz space, and the target is the closed ideal of MT-M24.
The open mapping theorem now proves continuity of the inverse in
these specified smooth topologies. Commutation with dilations is
immediate term by term. No zero-location result has been used.

Expanding the finite Euler product and regrouping the absolutely
convergent series gives (MA-TR-8) by the same divisor identity,
now restricted to square-free divisors composed of primes in $S$.
If $F$ is compactly supported, its corresponding even $f$ is
smooth and zero near the origin. Moreover
\[
 \int_{\mathbb R}f(t)\,dt
    =\int_0^\infty F(t)t^{-1/2}\,dt=\widehat F(1/2),
\]
so it belongs to $\mathscr S_{0,+}$. On any prescribed open
subinterval of $(a,b)$ the two functionals
$\widehat F(1/2)$ and $\widehat F(\lambda)$ are linearly independent:
otherwise the two smooth densities $t^{-1/2}$ and $t^{\lambda-1}$
would be proportional throughout that interval, forcing
$\lambda=1/2$, contrary to $\rho\ne1$. A linear map to
$\mathbb C^2$ with independent coordinate functionals is onto.
Choose $F$ with values $(0,1)$ under these two functionals.

Every $P_S\mathcal ZF$ belongs to the centered theta ideal, so its
Mellin value at $\lambda$ is zero by the full-jet description in
MT-M18/M24. This can also be seen first in $\Re z>1/2$ from
\[
 \widehat{\mathcal ZF}(z)=\zeta(z+1/2)\widehat F(z)
\]
and then by meromorphic continuation: the potential pole at $z=1/2$
is removed by the chosen moment. Cauchy--Schwarz on
$H_{\alpha,s}$ makes evaluation at $\lambda$ continuous with norm
at most the square root of the integral in (MA-TR-9). Applying it
to the difference proves that lower bound. Finally every $n>1$
in (MA-TR-8) is at least $q_S$; a nonzero summand requires $nt\le b$.
This proves the support assertion and exact eventual vanishing on
each compact interval.\hfill$\square$

\paragraph{Scope and input of MA-TR-101.}
The smooth isomorphism and Euler-tail identity are unconditional and
defined by the theta and M\"obius sums, without a spectral choice.
The lower bound is a conditional statement for each zero actually
retained by a chosen strip. A wide strip $\alpha>1/2$ retains every
nontrivial zero by the classical critical-strip theorem. For an
arbitrarily narrow strip, the existence of a critical-line zero already
suffices to instantiate the obstruction; this is classical, and is
strictly less than $A$. Such use of a known zero is recorded as a
zero-location input to the obstruction, never used to prove the
desired geometric equality. Alternatively the bound may simply be
left quantified over retained zeros, with no existence assertion.
The theorem does not prove that every other global-to-local
transport is impossible. It does rule out identifying compact
convergence of the finite Euler removals with convergence in the
specified Hilbert norm.

\subsection{An additional operator-order issue that is not formal}

\paragraph{MA-TR-102 (quotient trace class does not imply trace class after an orthogonal lift).}
There is a uniformly bounded strongly continuous representation
$V_t$ of $\mathbb R$ on a Hilbert space $H$, a closed invariant
subspace $J$, and a test $f\in C_c^\infty((0,\infty))$ such that
the integrated representation on $H/J$ is trace class, whereas
$C_fP_{J^\perp}$ on $H$ is not trace class. This remains true with
a simple purely imaginary quotient spectrum. Thus MT-M26/M28,
as quotient theorems, do not alone prove Masullo's separate
Theorem 10.4.2.8 about the opposite operator order.

\emph{Proof.}
Let $H=\ell^2(\mathbb N)\oplus\ell^2(\mathbb N)$ and
$J=\ell^2(\mathbb N)\oplus0$. Put
\[
 R_t=\operatorname{diag}(e^{int})_{n\ge1},\qquad
 K=\operatorname{diag}((in)^{-1})_{n\ge1},\qquad
 V_t=\begin{pmatrix}I&K(R_t-I)\\0&R_t\end{pmatrix}.
 \tag{MA-TR-10}
\]
The identity $K(R_{t+u}-I)=K(R_u-I)+K(R_t-I)R_u$ proves the
group law. Strong continuity follows from that of $R_t$, and
$\|K(R_t-I)\|\le2$ proves uniform boundedness. The quotient
action is $R_t$. For $F(z)=\int f(t)e^{zt}\,dt$ the integrated
matrix is
\[
 C_f=\begin{pmatrix}
 F(0)I&\operatorname{diag}((F(in)-F(0))/(in))\\
 0&\operatorname{diag}(F(in))
 \end{pmatrix}.
 \tag{MA-TR-11}
\]
Choose $f$ smooth, compactly supported away from zero, with
$F(0)\ne0$. Integration by parts gives rapid decay of $F(in)$,
so the quotient operator and $P_{J^\perp}C_f$ are trace class.
The upper-right corner of $C_fP_{J^\perp}$ has singular values
$|F(in)-F(0)|/n$, whose sum diverges. A corner of a trace-class
operator is trace class, so $C_fP_{J^\perp}$ is not.
The quotient eigenvalues $in$ are simple.\hfill$\square$

This is an abstract counterexample to an inference, not an arithmetic
counterexample to Theorem 10.4.2.8. The latter has its own unfinished
estimate on p.~464, (10.4.2.25)--(10.4.2.34). For the actual
arithmetic operators that additional estimate, or a valid replacement,
must be checked. F-SM-18 and F-SM-21 appropriately stated their
diagonal conclusions conditional on that actual lifted trace-class
input. The distinction is retained here.

\subsection{What the smooth transport does and does not accomplish}

The smooth map (MA-TR-7) explicitly contains the finite-prime
corrections. Its Mellin multiplier, wherever the Euler product
converges, is $1/\zeta(z+1/2)$, whose logarithmic derivative is
$-\zeta'/\zeta(z+1/2)$. With our translation convention the
corresponding distribution is
\[
 K(y)=\sum_{p,k\ge1}(\log p)p^{-k/2}\delta_{-k\log p}(y),
 \qquad
 (C_hK)(0)=\sum_{p,k\ge1}(\log p)p^{-k/2}h(p^k).
\]
Thus the centered contribution of the prime $p$ in this evaluation is
\[
 (\log p)\sum_{k\ge1}p^{-k/2}h(p^k)
\]
on tests supported above norm one. The reciprocal contribution is
obtained from the dual functional equation. This explains the
$\log2/\sqrt2$ correction in F-SM-21 using the actual arithmetic
map. It does not replace the missing trace estimate by a scalar
definition of a rank-one operator with the required trace.

For the source's central method, the outstanding assertion is that
the conjugated cutoff operators associated with $\mathcal J$ have
the correct common-space realization and a controlled remainder,
including the escape to norm zero in (MA-TR-9). The locally uniform
M\"obius limit is insufficient. An isomorphism of the smooth
geometric spaces supplies neither boundedness of its inverse in a
specified narrow Hilbert norm nor the convergence of nuclear traces.
The bounded Euler-module obstruction MT-M33 and the compatible
indentation obstruction MT-M37 continue to apply. They do not rule
out every non-equivariant family of cutoff operators.

In Grothendieck's terminology this is a question about transporting
an adelic trace on $C_{\mathbb Q}$ through a nonproper limiting
operation. Its Fredholm-determinant counterpart is the continuity
needed before taking $\det(1-wB)$ and its logarithmic derivative.
The exact archive comparison remains cote~1, scans~079--080 and
111--112, via ALPOGE\_BRIDGES BR-TR-01/02 and F-SM-11/12;
the archive pages themselves are not asserted freshly inspected.
At a finite spectral stage exterior powers and determinants live
in the usual $K_0$ and $\lambda$-ring formalism, but an analytic
limit is additional data. No theorem in EGA, FGA or SGA identifies
the above weighted Hilbert closure with a proper pushforward of
an algebraic sheaf; no such identification is claimed.
